\documentclass[11pt]{article}
\usepackage[margin=1in]{geometry}
\usepackage[T1]{fontenc}
\usepackage[utf8]{inputenc}
\usepackage{lmodern}
\usepackage{enumitem}

\title{Professor Presentation Cheat Sheet}
\author{Lunar}
\date{April 1, 2026}

\begin{document}
\maketitle
\thispagestyle{empty}

\section*{Opening}
Short update on the EarthRover project. Keep it simple: indoor, outdoor, marathon. For each one: what we chose, why we chose it, what worked, what failed.

\section*{Indoor}
\textbf{What we chose}
\begin{itemize}[leftmargin=*]
\item MBRA as the main indoor controller
\item exact checkpoint-step navigation
\item a cleaner runtime around it
\end{itemize}

\textbf{Why}
\begin{itemize}[leftmargin=*]
\item indoor is a known corridor problem
\item exact steps are cleaner than vague targets
\item MBRA worked better with simpler recovery
\end{itemize}

\textbf{Result}
\begin{itemize}[leftmargin=*]
\item reached \textbf{8 / 11} checkpoints
\item system became more stable and more understandable
\end{itemize}

\textbf{Main failures}
\begin{itemize}[leftmargin=*]
\item forward-only graph caused no-path issues
\item stale context could trap the robot
\item recovery still needed care
\end{itemize}

\section*{Outdoor}
\textbf{What we chose}
\begin{itemize}[leftmargin=*]
\item LogoNav as the main outdoor controller
\item OSM route expansion
\item safety layers around the runtime
\end{itemize}

\textbf{Why}
\begin{itemize}[leftmargin=*]
\item outdoor already had a working base runtime
\item the bigger need was safety and stability, not replacing the controller
\end{itemize}

\textbf{Result}
\begin{itemize}[leftmargin=*]
\item one run reached \textbf{all checkpoints}
\item the rest were around \textbf{50\%} success overall
\item \textbf{3 rounds failed}
\end{itemize}

\textbf{Main failures}
\begin{itemize}[leftmargin=*]
\item spinning
\item bad waypoint transitions
\item curbs and steps still weak
\end{itemize}

\section*{Marathon}
\textbf{What we chose}
\begin{itemize}[leftmargin=*]
\item same outdoor base runtime
\item stricter safety layers
\item focus on stability, not full autonomy
\end{itemize}

\textbf{What we added}
\begin{itemize}[leftmargin=*]
\item IMU safety
\item route corridor guard
\item rerouting from live GPS
\item better waypoint handling
\item clearer logs
\end{itemize}

\textbf{Result}
\begin{itemize}[leftmargin=*]
\item reached \textbf{1 checkpoint}
\item then the robot spun and toppled
\end{itemize}

\textbf{Reason}
\begin{itemize}[leftmargin=*]
\item likely unstable target handling after checkpoint transition
\item repeated alignment / turning in place
\item repeated aggressive turning made the robot unstable
\end{itemize}

\section*{Final Summary}
\begin{itemize}[leftmargin=*]
\item Indoor: best direction is MBRA-first, result \textbf{8 / 11}
\item Outdoor: best direction is LogoNav + route + safety layers, result \textbf{1 full success}, others around \textbf{50\%}
\item Marathon: main issue is stability after checkpoint and waypoint transitions
\end{itemize}

\section*{Numbers To Remember}
\begin{itemize}[leftmargin=*]
\item indoor: \textbf{8 / 11}
\item outdoor best run: \textbf{all checkpoints}
\item outdoor overall: \textbf{around 50\%}
\item outdoor failed rounds: \textbf{3}
\item marathon: \textbf{1 checkpoint, then toppled}
\end{itemize}

\section*{One Good Closing Line}
The project is much better understood now, and the remaining weaknesses are much clearer.

\end{document}
